Respuesta oral breve, sin tablas. Definimos la transformada para un parámetro real \(s\) mediante la integral impropia
\[
\mathcal L\{f\}(s)=\int_0^\infty e^{-st}f(t)\,dt.
\]
Sean \(f_1(t)=t^{-1/2}\) y \(f_2(t)=t^{-1}\) para \(0<t\le1\), ambas cero para \(t>1\), y \(f_1(0)=f_2(0)=0\). Sea además \(f_3(t)=e^{t^2}\). Decida para qué valores reales de \(s\) converge cada integral y justifique sin calcularla exactamente. Distinga lo que ocurre cerca de cero de lo que ocurre en el infinito. ¿Qué muestra el primer caso sobre las hipótesis \emph{suficientes} habituales de continuidad a trozos y orden exponencial? ¿Puede arreglarse el segundo cambiando sólo el valor de \(f_2(0)\)?
